\documentclass[11pt]{article}


\setlength{\parindent}{0pt}
\usepackage{xltxtra}
\usepackage{hyperref}
\hypersetup{hidelinks}
\usepackage{url}
\urlstyle{tt}
\usepackage{xcolor}
\definecolor{CVBlue}{RGB}{23,110,191}
\usepackage{calc}
\usepackage{graphicx}
\usepackage{tikz}
\usetikzlibrary{calc}
\usepackage{fontspec}
\usepackage{xeCJK}
\usepackage{enumitem}
\CJKsetecglue{} %% 取消中文与数字之间的间隙


%% 主文档字体设置
\setmainfont[
    Path = fonts/Main/,
    Extension = .otf,
    BoldFont = texgyretermes-bold.otf, % 加粗字体
]{texgyretermes-regular.otf} % 正文字体

% 中文字体设置
\setCJKmainfont[
    Path = fonts/hansans/,
    Extension = .ttf,
    BoldFont = NotoSansSC-Bold.ttf, % 加粗字体
]{NotoSansSC-Regular.otf} % 正文字体


\usepackage{relsize}
\usepackage{xspace}

% 使用 fontawesome（部分图标）
\usepackage{fontawesome} 

% A4纸，上下左右边距
\usepackage[
    a4paper,
    left=1.2cm,
    right=1.2cm,
    top=1.5cm,
    bottom=1cm,
    nohead
]{geometry}

\renewcommand{\baselinestretch}{1.5} % 行间距设为1.5

\usepackage{titlesec}
\usepackage{enumitem}
\setlist{noitemsep} % 取消列表项间的额外间距
%\setlist{nosep} % 取消所有垂直间距
\setlist[itemize]{topsep=0.25em, leftmargin=*}
\setlist[enumerate]{topsep=0.25em, leftmargin=*}

% --- 用于控制【不同项目之间】的垂直距离 ---
\newlength{\interProjectSpacing}
\setlength{\interProjectSpacing}{0.9em} % <--- 在此调整项目之间的距离
\newcommand{\projectsep}{\vspace{\interProjectSpacing}}

% --- 用于控制【项目标题】与下方【项目描述】的距离 ---
\newlength{\intraProjectTitleSep}
\setlength{\intraProjectTitleSep}{0.4em} % <--- 在此调整标题和描述的距离
\newcommand{\titlebreak}{\\[\intraProjectTitleSep]}

% --- 用于控制【项目描述】与下方【要点列表】的距离 ---
\newlength{\intraProjectListTopSep}
\setlength{\intraProjectListTopSep}{0.2em} % <--- 在此调整描述和列表的距离

% =======================================================================


\titleformat{\section}         % 定制 \section 命令 
{\large\bfseries\raggedright} % 将 section 标题设置为大号、粗体且左对齐
{}{0em}                      % 可用于为所有 section 添加前缀（如“章节...”）
{}                           % 可用于在标题前插入代码
[{\color{CVBlue}\titlerule}]  % 在标题后插入一条横线
\titlespacing*{\section}{0cm}{*1.6}{*1.2}



\begin{document}
\pagenumbering{gobble}

%%%% 利用tikz来定位照片
\begin{tikzpicture}[remember picture, overlay] 
    \node[anchor = north east] at ($(current page.north east)+(-2cm,-1.2cm)$) {\includegraphics[height=3cm]{avatar.jpg}};
  \end{tikzpicture}%
  %%%% 利用tikz来定位学校Logo，这里只在第一页显示
  \begin{tikzpicture}[remember picture, overlay] 
    \node[anchor = north west] at ($(current page.north west)+(0.5cm,+1.0cm)$) {\includegraphics[height=6cm]{zju.png}};
  \end{tikzpicture}%
\centerline{\LARGE\bfseries{王逸云}} 

\centerline{\normalsize{\faPhone\ 151-0617-6637 \quad \faEnvelopeO\ \href{mailto:3220102175@zju.edu.cn}{3220102175@zju.edu.cn}}} 
    
\section{\makebox[\widthof{\faGraduationCap}][c]{\color{CVBlue}\faGraduationCap}\ 教育背景}    
\textbf{浙江大学} \hfill 2022.9 -- 至今\\[0.5em] % 标题和正文间加一点距离
口腔医学（5+3）\quad 大四
学年绩点\titlebreak
\begin{itemize}[nosep]
    \item 2022 -- 2023: 4.42 17/226
    \item 2023 -- 2024: 4.50 5/50
    \item 2024 -- 2025: 4.47 6/50
    \item 2025 -- 至今: 4.40 -/50
\end{itemize}


\section{\makebox[\widthof{\faGraduationCap}][c]{\color{CVBlue}\faList}\ 获奖情况}

\textbf{浙江大学一等奖学金} \hfill 2024.12
\textbf{浙江省政府奖学金} \hfill 2023.12/2024.12/2025.12
浙江大学二等奖学金 \hfill 2023.12
浙江大学三等奖学金 \hfill 2025.12
浙江大学优秀学生 \hfill 2023.12/2024.12
浙江大学校级优秀团员 \hfill 2023.6/2024.6
浙江大学医学院优秀团干部 \hfill 2025.6
\textbf{全国大学生英语竞赛（NECCS）C类二等奖} \hfill 2023.6
浙江大学“学习贯彻二十大踔厉奋发新征程”寒假大学生社会实践活动优秀团队 \hfill 2023.4
浙江大学暑期社会实践“十佳团队”提名奖 \hfill 2023.12
浙江大学暑期社会实践优秀论文 \hfill 2023.12
浙江大学本科生优秀学长 \hfill 2024.6
浙江大学“三好杯”操舞比赛第五名 \hfill 2025.6

\section{\makebox[\widthof{\faUsers}][c]{\color{CVBlue}\faUsers}\ 学生工作}

\textbf{口腔（5+3）2202班班长} \hfill 2022.09 -- 至今 \titlebreak
工作描述：担任班级班长，统筹\textbf{日常管理与学风建设}，协调师生沟通，以 \textbf{责任心与服务意识}保障班级稳定有序运转。。
% 在 itemize 的选项中，使用 topsep=\intraProjectListTopSep 来控制上边距
\begin{itemize}[nosep, topsep=\intraProjectListTopSep]
    \item \textbf{沟通能力}：平均\textbf{每周}至少在班级群内发布\textbf{2条}通知信息，作为同学和老师的交流桥梁，高效传达信息并满足各方面需求。
    \item \textbf{可视化前端}：组织并执行 “跳脱出课本之外，徜徉于医海之中 —— 提升医学生沟通合作力与培育创新意识计划” \textbf{求是学院学生综合素质提升工程项目}，增强班级凝聚力、志愿服务意识与学习成绩。
\end{itemize}

% 使用 \projectsep 命令来分隔两个项目
\projectsep

% --- 第二个项目 ---
\textbf{医视野宣传中心编辑部干事 -- 部长} \hfill 2022.09 -- 2024.09 \titlebreak
工作描述：为学院宣传仁心仁术的品德风尚，为学子编写科研生活的时事资讯，在编辑部的工作中基于\textbf{公众号编写与排版}技术，独立完成内容创作并发表多篇原创文章。
\begin{itemize}[nosep, topsep=\intraProjectListTopSep]
    \item \textbf{宣传能力}：拥有两年新闻撰写经验，具备扎实的\textbf{现场记录} 能力及 \textbf{采访} 能力，可完成完整新闻采编工作。
    \item \textbf{写作水平}：全面掌握了从\textbf{录音记录、思维导图制作到Latex}的写作全链路技术。
    \item \textbf{可视化结果}：多篇新闻作品累计浏览量达 \textbf{7859} 内容受到医学院师生广泛关注与认可。
\end{itemize}

\projectsep

\textbf{口腔医学本科生第一党支部正式党员} \hfill 2025.12 -- 至今 \titlebreak

\section{\makebox[\widthof{\faCogs}][c]{\color{CVBlue}\faCogs}\ 科研情况}
    \textbf{SRTP校级项目：} γ分泌酶酶切LRP5形成LICD的作用机制研究 \hfill 2024.05 -- 2025.05\titlebreak
    \begin{itemize}[nosep, topsep=\intraProjectListTopSep]
        \item \textbf{实验操作}：深度学习 \textbf{Western Blot} 操作，成功验证揭示LRP5胞外段响应细胞外信号暴露γ分泌酶作用的GV结构域，形成游离LICD通过NLS核转位机制进入细胞核。
        \item \textbf{科研思维}：运用 \textbf{Pubmed} 精准筛选并深度阅读相关研究论文，形成系统的文献梳理与分析能力。
    \end{itemize}
    \textbf{生理科学实验结课项目：} 糖尿病对草酸钙肾结石严重程度的影响 \titlebreak
    \begin{itemize}[nosep, topsep=\intraProjectListTopSep]
        \item \textbf{实验操作}：熟练掌握 \textbf{Olympus数字切片扫描仪VS200} 操作，利用 \textbf{Image J软件} 计算视野下阳性区域面积占视野总面积的百分比以表示草酸钙结晶肾内沉积的相对水平。
    \end{itemize}
    
\section{\makebox[\widthof{\faInfo}][c]{\color{CVBlue}\faInfo}\ 其他}
\begin{itemize}[parsep=0.5ex] 
    \item \textbf{英语水平：} CET-4：655, CET-6：639
    \item \textbf{志愿活动：} 浙江大学四星级志愿者
\end{itemize}
\end{document}
